\include{LectureSlides/includes/metropolis_preamble}


% Title page info
\title[Introduction]{Introduction to Game Theory}
\author[ECON 5001]{ECON 5001\\P.J. Healy} \color{metrop}
\institute[OSU]{}
\date[]{\vfill {\tiny Updated \today\ at\ \DTMcurrenttime}}

\begin{document}

\frame{\maketitle}

\begin{frame}{What is Game Theory?}
\textit{Game theory} is the study of strategic interaction: settings where
each player's best choice depends on what the others do.
\br
\begin{itemize}
    \item Players, strategies, and payoffs
    \item Solution concepts: dominance, Nash equilibrium, and refinements
    \item Applications across economics, political science, biology, and CS
\end{itemize}
\end{frame}

\begin{frame}{A First Example: The Prisoner's Dilemma}
Two players simultaneously choose to \emph{Cooperate} or \emph{Defect}:
\br
\begin{center}
\begin{game}{2}{2}
      & $C$        & $D$        \\
$C$   & $-1,-1$    & $-3,\;0$   \\
$D$   & $0,-3$     & $-2,-2$
\end{game}
\end{center}
\br
\emph{Defect} strictly dominates \emph{Cooperate} for both players, so the
unique Nash equilibrium is $(D,D)$ --- even though both prefer $(C,C)$.
\end{frame}

\begin{frame}{Roadmap}
\begin{enumerate}
    \item Strategic-form games and dominance
    \item Nash equilibrium
    \item Mixed strategies
    \item Extensive-form games and backward induction
    \item Repeated games and reputation
\end{enumerate}
\end{frame}

\end{document}
